\section{方法选择：基础盈利估值与未来期权分层}

德赛西威已经具备稳定盈利和分部收入披露，因而前瞻P/E是基础业务的主方法。它将座舱的已确认盈利基础和智驾的增长质量放在同一个合并分母内。次要检查是产品线收入/毛利桥、公开EPS区间、现价隐含P/E以及市值与股本算术。我们不使用合并PS或PEG作为主方法，因为细分业务的客户/车型/ASP/增量费用尚不足以证明整家公司是纯高成长软件或纯智驾标的。Physical AI不进入2026E基础EPS，但以FY2028商业化结果的风险调整现值作为独立期权层；其敏感性和升级门槛见上一章。

\begin{exhibitbox}[价格、股本和估值算术]
\small
\begin{tabularx}{\textwidth}{L{3.2cm}Y Y}
\toprule
项目 & 计算 & 结果 \\
\midrule
市场价格 & 2026-07-23盘中快照 & CNY83.48 \\
总股本 & 2026Q1披露股本 & 5.96809294亿股 \\
市值 & CNY83.48 $\times$ 5.96809294亿股 & CNY49.822bn \\
基准EPS & CNY2.80bn / 5.96809294亿股 & CNY4.6916 \\
基础价值 & CNY4.6916 $\times$ 18.5x & CNY86.795，显示为CNY86.8 \\
基础业务多锚价值 & 0.8 $\times$ CNY86.8 + 0.2 $\times$ CNY85.0 & CNY86.44，显示为CNY86.4 \\
Physical AI风险调整期权 & FY2028结果按12\%折现；相互排斥情景加权 & CNY2.49/股 \\
合并参考值 & CNY86.44 + CNY2.49 & CNY88.93，显示为CNY88.9 \\
上行空间 & CNY88.9 / CNY83.48 - 1 & +6.50\% \\
\bottomrule
\end{tabularx}
\sourcenote{市场数据包、本院模型；显示值存在四舍五入。}
\end{exhibitbox}

算术的透明度很重要：它使读者可以不同意收入、利润、倍数或期权概率假设，却仍能复算合并参考值。基准没有使用单季度年化、未披露订单、客户销量倒推或过时券商目标，也没有把AI Cube、机器人域控、低速无人车、Thor POC或Robotaxi合作加入2026E EPS。任何改变基础倍数的理由都必须与下一季收入、毛利、现金或客户链证据相连；任何新增Physical AI价值必须先跨过独立的商业化门槛。

\section{Physical AI：风险调整期权进入合并参考值}

这不是把远期故事并入2026E盈利，也不是把主题热度变成倍数溢价。我们把FY2028可能结果折现两年至2026年7月23日，再以相互排斥的本院风险权重加总。55\%的“未形成可衡量商业转化”权重反映公司仍未披露具名客户、实际交付、收入、毛利或现金闭环；35\%的“首次商业化规模”权重反映已有产品、定点和2026计划量产；只有10\%给予“可复制平台规模”，因为多客户、多场景和连续回款尚无披露。概率是本院承保假设，不是公司指引或一致预期。

\begin{exhibitbox}[FY2028 Physical AI真实期权：概率加权现值]
\small
\begin{tabularx}{\textwidth}{L{2.55cm}Y Y Y Y}
\toprule
FY2028结果 & 概率 & 条件经济学与估值 & 折现至2026年 & 概率加权价值 \\
\midrule
未形成可衡量商业转化 & 55\% & 无可审计收入/利润闭环 & CNY0.00 & CNY0.00 \\
首次商业化规模 & 35\% & 收入CNY1.0bn、净利率8\%、22x P/E；条件价值CNY2.95/股 & CNY2.35/股 & CNY0.82/股 \\
可复制平台规模 & 10\% & 收入CNY5.0bn、净利率10\%、25x P/E；条件价值CNY20.94/股 & CNY16.69/股 & CNY1.67/股 \\
\midrule
\textbf{风险调整期权价值} & \textbf{100\%} & FY2028结果按12\%折现两年 & -- & \textbf{CNY2.49/股} \\
\bottomrule
\end{tabularx}
\sourcenote{本院FY2028条件情景与风险权重；股本5.968亿股。收入、利润率、P/E、概率和折现率均非公司指引。}
\end{exhibitbox}

因此，CNY86.4是可验证汽车电子业务的多锚价值，CNY2.49是Physical AI的低置信度、风险调整现值，二者相加形成CNY88.9的合并参考值。CNY2.49并不表示公司已获得等额确定价值：在可复制平台规模分支中，未经概率调整的现值为CNY16.69/股，但现时仅以10\%计入。后续商业化事实应改变概率和条件经济学，而不是用更高主业P/E替代验证。

\section{三情景与预期桥}

熊市CNY69.1对应收入CNY35.5bn、归母CNY2.50bn、EPS CNY4.19和16.5倍P/E。它不是灾难情景，而是收入放缓、智能驾驶毛利继续受压及营运资本占用抬升时的正常估值收缩。基准CNY86.8要求收入CNY38.0bn、归母CNY2.80bn、EPS CNY4.69和18.5倍；牛市CNY105.6要求CNY40.0bn、CNY3.15bn、CNY5.28及20倍，前提是量产、产品结构和利润转化有明确验证。

\begin{exhibitbox}[预期桥：市场在为什么付费]
\small
\begin{tabularx}{\textwidth}{L{2.7cm}Y Y}
\toprule
驱动 & 市场/公开资料能够支持的内容 & 仍需验证的内容 \\
\midrule
收入增长 & 2025已增长，智驾增速快；公开2026E收入约CNY38.4bn & 具体客户、车型、SOP与订单确认节奏 \\
利润率 & 座舱/智驾已有分部毛利率 & 智驾成本、SoC/BOM、项目价格与增量费用 \\
倍数 & 盘中约17.1x公开2026E EPS & 20x牛市倍数需要的增长持续期与质量 \\
产品上移 & 新平台、HUD、区域控制器、中央计算有进展 & 客户实名、ASP、收入、毛利和现金贡献 \\
\bottomrule
\end{tabularx}
\sourcenote{公司年报、券商原始报告、市场快照。}
\end{exhibitbox}

当前价格接近基础业务价值的原因是市场已给予成长预期。基础目标的微小上行来自经营兑现而非强行提高倍数；叠加低置信度期权后，合并上行约6.5\%，仍不足以改变“观察而非追高型核心配置”的行动纪律。但这不能替代对未来机会的研究：若Physical AI跨过客户、量产、收入、毛利与现金门槛，须更新单列收入/EPS桥和概率，而不是在此处直接上调主业P/E。

\section{多锚目标与行动逻辑}

基本面锚CNY86.8占80\%，因为它来自可复算的收入、利润、股本和倍数；市场锚CNY85.0占20\%，因为市场价格和公开EPS反映可观察的成长定价；Street价格锚为0\%，因为并无当前可审计目标价/方法。权重相加为100\%。

这个结构避免三种常见错误：第一，把自己的保守P/E机械地压过所有市场信息，造成不解释情绪溢价的“卖出”；第二，把无来源的券商目标或搜索摘要变成支持性锚；第三，把尚无财务闭环的Physical AI叙事直接变成主业倍数溢价。若中报验证了收入和毛利，本院会先提高基础盈利的置信度，再决定是否提高市场或基本面锚；若机器人/无人车出现商业化闭环，则更新已单列的概率/情景模型。

\begin{riskbox}[估值失效条件]
基础目标失效于收入持续低增、智能驾驶毛利继续显著下滑、应收/库存/减值使现金转换恶化，或客户/芯片/海外风险转化为实际交付偏差。牛市价值仅在新增平台闭合客户--SOP--收入/毛利链条后才可讨论。
\end{riskbox}
